\chapter{Project Structure from GitHub}

After cloning the AUTpaTT v3 repository from GitHub, the expected directory structure is:

\begin{itemize}
  \item \textbf{/firmware/} -- STM32 Nucleo-L010RB source code, configuration files, and precompiled binaries.
  \item \textbf{/app/} -- Python app (using PyQt6) graphical user interface for system control and visualization.
  \item \textbf{/mechanical/} -- 3D models (STL, CAD), mechanical drawings, and assembly diagrams.
  \item \textbf{/electronics/} -- Schematics, wiring diagrams, and optional PCB design files.
  \item \textbf{/docs/} -- Additional documentation, figures, and exported PDFs.
  \item \textbf{/config/} or root file \texttt{params.json} -- System configuration parameters.
\end{itemize}

Each folder should include a minimal \texttt{README} that describes its content and how it is used.

You can add or update this structure as the project evolves.
